\documentclass[../../../../main.tex]{subfiles}
\begin{document}

各世代ごとに，各個体が，他の個体とは独立に，
確率 $p$ で1個，確率 $1-p$ で2個の新しい個体を次の世代に残し，
それ自身は消滅する細胞がある．
いま，第0世代に1個であった細胞が，
第 $n$ 世代に $m$ 個となる確率を，
$P_n(m)$ とかくことにしよう．

\begin{figure}[htbp]
  \centering
  \begin{tikzpicture}[scale=1.2, >=stealth]
    \node[circle, draw, inner sep=2pt, label=above:第0世代] (G0) at (0,0) {};

    \node[circle, draw, inner sep=2pt, label=above:第1世代] (G1a) at (1.5, 0.8) {};
    \node[circle, draw, inner sep=2pt, label=below:第1世代] (G1b) at (1.5, -0.8) {};
    \draw (G0) -- (G1a);
    \draw (G0) -- (G1b);

    \node[circle, draw, inner sep=2pt, label=above:第2世代] (G2a) at (3.0, 1.4) {};
    \node[circle, draw, inner sep=2pt, label=above:第2世代] (G2b) at (3.0, -0.3) {};
    \node[circle, draw, inner sep=2pt, label=below:第2世代] (G2c) at (3.0, -1.3) {};

    \draw (G1a) -- (G2a);
    \draw (G1b) -- (G2b);
    \draw (G1b) -- (G2c);

    \node[circle, draw, inner sep=2pt, label=above:第3世代] (G3a) at (4.5, 1.6) {};
    \node[circle, draw, inner sep=2pt, label=above:第3世代] (G3b) at (4.5, 0.0) {};
    \node[circle, draw, inner sep=2pt, label=above:第3世代] (G3c) at (4.5, -1.0) {};
    \node[circle, draw, inner sep=2pt, label=below:第3世代] (G3d) at (4.5, -1.8) {};

    \draw (G2a) -- (G3a);
    \draw (G2b) -- (G3b);
    \draw (G2c) -- (G3c);
    \draw (G2c) -- (G3d);
  \end{tikzpicture}
\end{figure}

$n$ を自然数とするとき，$P_n(1)$，$P_n(2)$，$P_n(3)$ を求めよ．

\end{document}
